\section{Representing Boolean Functions}

A Boolean function is a function that takes one or more binary variables, i.e 0's and 1's as input, and produces a single binary output (0 or 1).

If a function has (n) variables, then the number of possible input combinations is:
\[
2^n
\]

Representing a Boolean function means expressing the function in a form that clearly shows the relationship between its inputs and output.

There are just a couple of ways of representing boolean functions, some of these main ways of representing Boolean functions are:
\begin{itemize}
    \item Truth Tables
    \item Sum of Products (SOP)
    \item Product of Sums (POS)
\end{itemize}

\subsection{Truth Tables}

A truth table lists all possible combinations of input values and their corresponding outputs.

\subsubsection*{Example 1}
\[
f(x,y) = x + y
\]

\begin{center}
\begin{tabular}{|c|c|c|}
\hline
x & y & f \\
\hline
0 & 0 & 0 \\
0 & 1 & 1 \\
1 & 0 & 1 \\
1 & 1 & 1 \\
\hline
\end{tabular}
\end{center}

\subsubsection*{Example 2}
\[
f(x,y) = xy
\]

\begin{center}
\begin{tabular}{|c|c|c|}
\hline
x & y & f \\
\hline
0 & 0 & 0 \\
0 & 1 & 0 \\
1 & 0 & 0 \\
1 & 1 & 1 \\
\hline
\end{tabular}
\end{center}

\subsubsection*{Example 3}
\[
f(x,y,z) = x + yz
\]

\begin{center}
\begin{tabular}{|c|c|c|c|c|}
\hline
x & y & z & yz & f \\
\hline
0 & 0 & 0 & 0 & 0 \\
0 & 0 & 1 & 0 & 0 \\
0 & 1 & 0 & 0 & 0 \\
0 & 1 & 1 & 1 & 1 \\
1 & 0 & 0 & 0 & 1 \\
1 & 0 & 1 & 0 & 1 \\
1 & 1 & 0 & 0 & 1 \\
1 & 1 & 1 & 1 & 1 \\
\hline
\end{tabular}
\end{center}

\subsection{Sum of Products (SOP)}

A Boolean function is in Sum of Products form if it is written as a product(AND) of sum(OR) terms.

\subsubsection*{Minterm Definition}
A minterm is a product term in which each variable appears exactly once, either complemented or uncomplemented, and corresponds to a row where the function value is 1.

\subsubsection*{Canonical SOP}
A canonical SOP is the sum of all minterms where the function equals 1.

\subsubsection*{Example 1}

\begin{center}
\begin{tabular}{|c|c|c|}
\hline
x & y & f \\
\hline
0 & 0 & 1 \\
0 & 1 & 0 \\
1 & 0 & 1 \\
1 & 1 & 0 \\
\hline
\end{tabular}
\end{center}

\[
f = x'y' + xy' = \sum m(0,2)
\]

\subsubsection*{Example 2}

\[
f = x'y + xy' + xy = \sum m(1,2,3)
\]

\subsubsection*{Example 3}

\[
f = x'y'z + x'yz + xy'z' + xyz' = \sum m(1,3,4,6)
\]

\subsubsection*{Maxterm Definition}
A maxterm is a sum term in which each variable appears exactly once and corresponds to a row where the function value is 0.

\subsubsection*{Canonical POS}
A canonical POS is the product of all maxterms where the function equals 0.

\subsubsection*{Example 1}

\[
f = (x + y')(x' + y') = \prod M(1,3)
\]

\subsubsection*{Example 2}

\[
f = (x + y)(x' + y) = \prod M(0,2)
\]

\subsubsection*{Example 3}

\[
f = (x + y + z)(x + y' + z')(x' + y + z')(x' + y' + z)
\]
\[
= \prod M(0,3,5,6)
\]

\subsection{Function Completeness}

A set of Boolean operations is functionally complete if every Boolean function can be expressed using only operations from that set.

\subsubsection*{Standard Complete Set}
\[
\{AND, OR, NOT\}
\]

\subsubsection*{Example 1}

\[
f(x,y) = x'y + xy'
\]

\subsubsection*{NAND Completeness}

\[
x' = x \uparrow x
\]
\[
xy = (x \uparrow y)'
\]

\subsubsection*{NOR Completeness}

\[
x' = x \downarrow x
\]

\subsection{Not Functionally Complete Sets}
\begin{itemize}
    \item \{AND, OR\}
    \item \{AND, NOT\}
    \item \{OR, NOT\}
\end{itemize}

\subsection{Conclusion}

Boolean functions can be represented using truth tables, SOP, and POS forms. Truth tables provide a complete description, while SOP and POS provide structured algebraic forms. Function completeness ensures that all Boolean functions can be constructed from a given set of operations.

